% 原书第 209—212 页；第 13.6 节导言与第 13.6.1 小节全文。
\section{小样本模拟}
\label{sec:ch13-small-sample-simulations}

下面的模拟研究旨在展示：（1）所提方法的实施；（2）通过数据自适应估计条件结局均值，可能获得的精度与功效改善；（3）通过协同估计暴露机制，进一步获得的精度与功效改善；（4）维持置信区间的名义覆盖概率。所有模拟都在 R 中进行（R Development Core Team，2016）。完整 R 代码见 Balzer et al.（2016b）。

\subsection{研究 1}
\label{subsec:ch13-study1}

对于每个单位 $i=\{1,\ldots,n\}$，从均值为 $0$、方差为 $1$ 的多元正态分布中抽样，生成九个基线协变量。前三个协变量 $\{W1,W2,W3\}$ 内部，以及第二组三个协变量 $\{W4,W5,W6\}$ 内部的相关系数为 $0.5$；其余协变量 $\{W7,W8,W9\}$ 之间的相关系数为 $0$。暴露 $A$ 以保证总体处理分配平衡的方式随机化。在无匹配试验中，将干预随机分配给 $n/2$ 个单位，其余 $n/2$ 个单位作为对照。在配对匹配试验中，采用非二分匹配算法 \texttt{nbpMatch}，按协变量 $\{W1,\ldots,W6\}$ 对单位配对（Beck et al.，2016），再在所得匹配对内随机分配暴露 $A$。回顾一下，$A$ 是二值指示变量，单位被分配至处理组时等于 $1$，被分配至对照组时等于 $0$。随后，为每个单位生成结局
\[
 Y=0.4A+0.25(W1+W2+W4+W5+U_Y)+0.25A(W1+U_Y),
\]
其中，$U_Y$ 抽自标准正态分布。还通过干预将 $A$ 设为 $a$，生成反事实结局 $Y_1$ 和 $Y_0$。为了反映早期临床试验和 CRT 中常见的有限样本量，选择 $n=40$。因此，配对匹配试验中有 $n/2=20$ 个条件独立单位。

对于每种研究设计，即无匹配或有匹配，将这一数据生成过程重复 $2500$ 次。回顾一下，样本效应，即式 (13.2)，是一个数据自适应参数；每次重新选择单位，它的值都会变化。因此，每次重复时，都将反事实结局之差的样本均值作为 SATE。SATE 的范围为 $0.22$—$0.59$，均值为 $0.40$。相比之下，总体效应，即式 (13.1)，是一个常数；它通过对包含 $900{,}000$ 个单位的总体中的反事实结局之差取平均来计算。PATE 的真值为 $0.40$。

我们将未经调整的估计量与采用不同协变量调整方法的 TMLE 比较。具体而言，实施这样一种 TMLE 算法：使用包含干预 $A$ 和无关协变量 $W9$ 主效应的线性工作模型，初始估计条件结局均值 $\bar Q_0(A,W)$，同时将暴露机制视为已知，即 $g_0(A\mid W)=0.5$。这种方法等价于标准 MLE，代表一种不理想的情形：研究者预先指定调整的协变量，不能预测结局。

还实施了在步骤 1 中，用数据自适应方法初始估计条件结局均值的 TMLE。算法库包含十个线性回归工作模型，每个模型都有截距、暴露 $A$ 的主效应，以及一个基线协变量的主效应。候选协变量为 $\{\varnothing,W1,\ldots,W9\}$，其中 $\varnothing$ 对应于未经调整的估计量。表现度量，即风险函数，是与目标参数和研究设计相适应的 TMLE 估计渐近方差。在无匹配试验中，基于留一交叉验证选择候选工作模型；在有匹配试验中，则基于留一对交叉验证。

我们也实施了 C-TMLE，将步骤 1 中初始估计条件结局均值的数据自适应方法，与步骤 2 中靶向更新的数据自适应方法结合。对于后者，用于估计暴露机制的候选算法库包含十个 logistic 回归工作模型，每个模型都有截距及一个基线协变量的主效应，候选协变量为 $\{\varnothing,W1,\ldots,W9\}$。C-TMLE 使用相同的损失函数和交叉验证方案。

对于未经调整的估计量和 MLE，推断基于估计的影响曲线。对于数据自适应 TMLE，推断基于影响曲线的交叉验证估计，见第 13.5 节。我们假设，标准化估计量服从 Student $t$ 分布：无匹配试验的自由度为 $n-2=38$，有匹配试验的自由度为 $n/2-1=19$。

\noindent\textbf{结果。}
表 13.1 展示了估计量在 $2500$ 个模拟数据集上的表现。具体列出了均方误差（MSE）、相对均方误差（rMSE）、平均标准误估计 $\hat\sigma$、实际达到的功效，以及 $95\%$ 置信区间的覆盖率。正如预期，匹配提高了效率。例如，无匹配试验中，未经调整估计量的 MSE 超过配对匹配试验中的两倍。此外，对于配对匹配试验，将样本效应而非总体效应作为目标，会显著提高实际达到的功效：同一个未经调整的估计量，以 PATE 为目标时功效为 $36\%$，以 SATE 为目标时为 $53\%$。对于无匹配试验，以样本参数为目标提高了效率，即 MSE 更小；但由于 SATE 的方差估计量偏于保守，这并未直接转化为更高的功效。

\begin{table}[htbp]
 \centering
 \caption{模拟研究 1 中估计量表现的汇总}
 \label{tab:ch13-study1}
 \begingroup
 \small
 \setlength{\tabcolsep}{3pt}
 \renewcommand{\arraystretch}{1.16}
 \begin{tabular}{lrrrrrrrrrr}
  \toprule
  &\multicolumn{5}{c}{PATE}&\multicolumn{5}{c}{SATE}\\
  \cmidrule(lr){2-6}\cmidrule(lr){7-11}
  &MSE&rMSE&$\hat\sigma$&功效&覆盖率&MSE&rMSE&$\hat\sigma$&功效&覆盖率\\
  \midrule
  \multicolumn{11}{l}{\textbf{无匹配}}\\
  未调整&$6.8\mathrm E{-2}$&1.00&0.25&0.34&0.94&$6.4\mathrm E{-2}$&1.06&0.25&0.34&0.94\\
  MLE&$6.9\mathrm E{-2}$&0.98&0.25&0.35&0.94&$6.5\mathrm E{-2}$&1.04&0.25&0.35&0.94\\
  TMLE&$4.5\mathrm E{-2}$&1.49&0.20&0.48&0.94&$4.2\mathrm E{-2}$&1.62&0.20&0.48&0.95\\
  C-TMLE&$4.3\mathrm E{-2}$&1.57&0.20&0.48&0.95&$4.0\mathrm E{-2}$&1.70&0.20&0.48&0.96\\
  \addlinespace[.3em]
  \multicolumn{11}{l}{\textbf{有匹配}}\\
  未调整&$3.2\mathrm E{-2}$&2.10&0.22&0.36&0.99&$2.9\mathrm E{-2}$&2.31&0.18&0.53&0.97\\
  MLE&$3.4\mathrm E{-2}$&2.01&0.22&0.37&0.98&$3.1\mathrm E{-2}$&2.19&0.18&0.53&0.96\\
  TMLE&$2.6\mathrm E{-2}$&2.64&0.19&0.51&0.98&$2.3\mathrm E{-2}$&2.93&0.16&0.65&0.96\\
  C-TMLE&$2.5\mathrm E{-2}$&2.71&0.18&0.53&0.98&$2.2\mathrm E{-2}$&3.03&0.15&0.67&0.96\\
  \bottomrule
 \end{tabular}
 \endgroup
 \par\smallskip
 \begin{minipage}{\linewidth}\footnotesize
 各行表示研究设计与估计量：未经调整的估计量、调整 $W9$ 的 MLE、数据自适应选择初始估计量的 TMLE，以及采用数据自适应选择的 C-TMLE。各列表示估计量表现：MSE 为偏差平方加方差；rMSE 为无匹配试验中 PATE 的未经调整估计量的 MSE，除以另一个估计量的 MSE；$\hat\sigma$ 为平均标准误估计；其余为功效和覆盖率。
 \end{minipage}
\end{table}

在所有情形中，调整无关协变量 $W9$ 的 MLE，其 MSE 都比其他估计量更差。这说明了依赖单个预先指定调整变量的潜在风险。事实上，通过数据自适应选择 $\bar Q_0(A,W)$ 的初始估计量，TMLE 的精度优于未经调整的估计量和 MLE。协同估计暴露机制 $g_0(A\mid W)$，又进一步提高了精度。例如，考虑无匹配试验中 PATE 的估计。未经调整估计量的 MSE，分别是 TMLE 的 $1.49$ 倍和 C-TMLE 的 $1.57$ 倍；三者实际达到的功效分别为 $34\%$、$48\%$ 和 $48\%$。再如，考虑配对匹配试验中，检测 SATE 是否异于零所达到的功效。采用未经调整的估计量，或者采用调整无关协变量 $W9$ 的 MLE，功效均为 $53\%$。纳入用于初始估计 $\bar Q_0(A,W)$ 的交叉验证选择器后，TMLE 达到了 $65\%$ 的功效；再进一步纳入暴露机制 $g_0(A\mid W)$ 的协同估计后，C-TMLE 达到了 $67\%$ 的功效。

总体而言，在配对匹配试验中，以 SATE 为目标的 C-TMLE 获得了最高效率。事实上，无匹配试验中总体参数的未经调整估计量，其 MSE 是配对匹配试验中样本效应的 C-TMLE 的三倍。所有情形下，置信区间覆盖率都维持在接近或高于 $95\%$ 的名义水平。
\FloatBarrier

